\begin{textsample}{POS Dim 6 – vem\_ed\_35 – Score -7.00 – vem\_ed\_35/t191.txt}  \label{ex:f6_pos_154}
Artigo de Opinião Cooperação que conecta, transforma e amplia horizontes O tema do 5º IVES CGESG 2026, Cooperation in COIL Virtual Exchange: hurdles and bridges (Cooperação em Intercâmbios Virtuais do tipo COIL: obstáculos e pontes), convida a refletir sobre como a educação pode ultrapassar fronteiras e criar conexões significativas entre culturas, instituições e pessoas.
O simpósio teve como propósito discutir perspectivas sobre a Aprendizagem Colaborativa Internacional realizada on-line (Collaborative Online International Learning – COIL) e os Projetos Colaborativos Internacionais (PCIs) desenvolvidos nas Fatecs.
O IVES CGESG promove a integração entre instituições de ensino superior brasileiras e internacionais e contribui para a disseminação de boas práticas e para a formação docente em metodologias de Internacionalização em Casa, fortalecendo a educação conectada ao cenário global.
Os resultados mostram a relevância desse trabalho: entre 2013 e 2026, ocorreram mais de 800 PCIs nas Fatecs, envolvendo 59 \textbf{unidades}, 96 instituições internacionais e aproximadamente 32 mil estudantes do Brasil e do exterior.
Esses números revelam não apenas alcance institucional, mas principalmente a capacidade de transformar trajetórias acadêmicas por meio da colaboração internacional.
Quando abrimos novas formas de aprendizagem, descobrimos o mundo de maneira mais humana, ampla e transformadora.
Esse movimento nos permite enxergar além das nossas fronteiras e ampliar nossa visão de mundo, fortalecendo nossa comunicação e nosso desenvolvimento profissional como educadores.
A universidade do futuro será verdadeiramente significativa quando estiver preparada para dialogar com diferentes culturas e instituições formadoras de pessoas.
É nessa troca que surgem oportunidades valiosas de compartilhar saberes e realizar trabalhos acadêmicos que aproximam estudantes de diferentes realidades.
O impacto dessa experiência é profundo e transformador.
Por meio dos PCIs, milhares de estudantes vivenciam experiências internacionais e ampliam seus horizontes acadêmicos e pessoais sem precisar sair da sua cidade.
Isso torna a educação mais acessível, inclusiva e alinhada às necessidades do nosso tempo.
A neurociência também contribui como referência importante nesse processo, ao reforçar o protagonismo estudantil e a construção ativa do conhecimento.
Quando o estudante participa, interage e se reconhece como parte essencial da aprendizagem, o conhecimento ganha mais significado e se torna mais duradouro.
É verdade que o encontro entre diferentes culturas pode trazer desafios e, em alguns momentos, gerar conflitos.
No entanto, esses obstáculos também representam oportunidades valiosas de crescimento quando existe respeito, escuta ativa e compreensão sobre o contexto de cada participante e sobre com quem estamos construindo essa experiência de aprendizagem.
Por isso, fortalecer iniciativas como os PCIs é investir em uma educação mais humana, inclusiva e transformadora.
Erguer pontes entre culturas, conhecimentos e experiências é ampliar possibilidades e transformar vidas.
A educação tem o poder de conectar pessoas e abrir caminhos antes inimagináveis.
Quando escolhemos aprender com o outro, crescemos juntos rumo a uma educação inovadora e verdadeiramente global.
Esse trabalho ganha ainda mais significado por meio do PCI desenvolvido entre as Fatecs Pompeia, Zona Leste e Guarulhos e a Bukidnon State University, nas Filipinas.
Docentes e estudantes das duas instituições conectam culturas e realidades distintas em uma experiência enriquecedora de troca de conhecimentos, vivências e perspectivas.
Essa colaboração fortalece o protagonismo estudantil e mostra como a educação pode ultrapassar fronteiras geográficas para aproximar pessoas, ampliar horizontes e promover aprendizagens significativas.
Esse intercâmbio de experiências contribui diretamente para a formação de estudantes mais preparados para atuar em um mundo diverso, global e em constante transformação.
É por meio de iniciativas como essa que percebemos o verdadeiro valor da internacionalização no ensino: não apenas como uma oportunidade acadêmica, mas como uma experiência humana transformadora.

% matched lemmas: unidade
\end{textsample}
